\chapter{System Identification of Hexapole Magnetic Flux Dynamics}

\section{Chapter Objective and Evidence Contract}

This chapter develops the system-identification stage of the thesis workflow. The objective is to convert open-loop observations into model information that is directly usable for control design and downstream force-generation analysis. The evidence contract of this chapter is to present simulation and experimental results in the same thread, so model quality is assessed by control relevance rather than by fitting score alone.

\section{Open-Loop Calibration Context and Data Interface}

This section defines the calibration context used in this dissertation, including excitation protocol, measurement channels, operating conditions, and data acquisition interface. The purpose is to make the identification input reproducible and to clarify what uncertainty sources are present before controller design begins.

\section{Frequency-Response Construction and Observable Phenomena}

This section describes how frequency-domain information is constructed from calibration data and how key physical phenomena are interpreted at this stage, including coupling behavior, remanence- and hysteresis-related effects, and frequency-dependent response variation. The focus is on observable evidence and its relevance to model structure selection.

\section{Control-Use Model Structure and Reduction Logic}

This section presents the model representation used for control development and the reduction logic that maps detailed channel behavior into a controller-usable form. The objective is to preserve physically meaningful coupling information while keeping a practical interface for the control-design stage.

\section{Fitting Strategy and Control-Relevant Quality Criteria}

This section defines the fitting strategy and quality criteria used in identification, with emphasis on low-frequency control relevance and predictable downstream behavior. The criteria are organized to support later design decisions in flux control and force generation.

\section{Simulation and Experimental Validation}

This section provides paired simulation and experimental validation for the identification stage. The reported results verify consistency between modeled behavior and measured behavior under tested operating conditions, and establish the handoff artifacts required by Chapter 4.

\section{Chapter Summary and Interface to Chapter 4}

This chapter concludes by summarizing what model information is accepted for control design, what uncertainty remains, and how identification outputs are passed to the flux-control stage. This closes the first link of the ID-to-Flux-to-Force evidence chain.
